\chapter{常用公式}
\label{app:formulas}

\begingroup
\fontsize{12pt}{20pt}\selectfont
\setlength{\parskip}{4pt}
\setlength{\abovedisplayskip}{7pt}
\setlength{\belowdisplayskip}{7pt}
\setlength{\abovedisplayshortskip}{4pt}
\setlength{\belowdisplayshortskip}{4pt}
\setlength{\LTpre}{6pt}
\setlength{\LTpost}{6pt}
\setlength{\jot}{1pt}

以下公式按对象集中列出，具体结论的定义、证明与使用见正文。求导、积分及三角函数公式中的角度均采用弧度制。

\section{数列与极限}
\begin{equation*}
 \begin{aligned}
 &a_n=a_1+(n-1)d,&
 &S_n=\frac{n(a_1+a_n)}2,\\
 &a_n=a_1q^{n-1},&
 &S_n=\frac{a_1(1-q^n)}{1-q}\quad(q\ne1).
 \end{aligned}
\end{equation*}
等比数列\(q=1\)时\(S_n=na_1\)；\(|q|<1\)时\(\sum_{n=1}^\infty a_1q^{n-1}=a_1/(1-q)\)。由部分和恢复通项的关系为\(a_n=S_n-S_{n-1}\)，其中\(S_0=0\)。
\begin{equation*}
 \begin{aligned}
 \sum_{k=1}^n k&=\frac{n(n+1)}2,&
 \sum_{k=1}^n k^2&=\frac{n(n+1)(2n+1)}6,\\
 \sum_{k=1}^n k^3&=\left(\frac{n(n+1)}2\right)^2,&
 e&=\lim_{n\to\infty}\left(1+\frac1n\right)^n.
 \end{aligned}
\end{equation*}
基本函数极限有\(\lim_{x\to0}\sin x/x=1\)、\(\lim_{x\to0}(e^x-1)/x=1\)、\(\lim_{x\to0}\ln(1+x)/x=1\)。

\section{三角恒等式}
\begin{equation*}
 \begin{aligned}
 \sin(\alpha\pm\beta)&=\sin\alpha\cos\beta\pm\cos\alpha\sin\beta,\\
 \cos(\alpha\pm\beta)&=\cos\alpha\cos\beta\mp\sin\alpha\sin\beta,\\
 \tan(\alpha\pm\beta)&=\frac{\tan\alpha\pm\tan\beta}{1\mp\tan\alpha\tan\beta},\\
 \sin2\alpha&=2\sin\alpha\cos\alpha,\\
 \cos2\alpha&=\cos^2\alpha-\sin^2\alpha=1-2\sin^2\alpha=2\cos^2\alpha-1,\\
 \sin^2\alpha&=\frac{1-\cos2\alpha}2,\qquad
 \cos^2\alpha=\frac{1+\cos2\alpha}2.
 \end{aligned}
\end{equation*}
记\(m=(\alpha+\beta)/2\)、\(d=(\alpha-\beta)/2\)，则
\begin{equation*}
 \begin{aligned}
 \sin\alpha+\sin\beta&=2\sin m\cos d,&
 \sin\alpha-\sin\beta&=2\cos m\sin d,\\
 \cos\alpha+\cos\beta&=2\cos m\cos d,&
 \cos\alpha-\cos\beta&=-2\sin m\sin d.
 \end{aligned}
\end{equation*}
设\(R=\sqrt{a^2+b^2}>0\)，取\(\cos\varphi=a/R\)、\(\sin\varphi=b/R\)，则\(a\sin x+b\cos x=R\sin(x+\varphi)\)。三角形中，若边\(a,b,c\)分别对角\(A,B,C\)，外接圆半径为\(R_\triangle\)，则
\begin{equation*}
 \frac a{\sin A}=\frac b{\sin B}=\frac c{\sin C}=2R_\triangle,\qquad
 c^2=a^2+b^2-2ab\cos C,\qquad S=\frac12ab\sin C.
\end{equation*}

\section{向量与矩阵}

在正交单位基下，\(u\cdot v=\sum_i u_iv_i\)，\(\|u\|=\sqrt{\sum_i u_i^2}\)。非零向量满足
\begin{equation*}
 \cos\theta=\frac{u\cdot v}{\|u\|\|v\|},\qquad
 \operatorname{proj}_v u=\frac{u\cdot v}{v\cdot v}v.
\end{equation*}
在三维右手正交单位基下，叉积的坐标为
\begin{equation*}
 u\times v=(u_2v_3-u_3v_2,\ u_3v_1-u_1v_3,\ u_1v_2-u_2v_1).
\end{equation*}
平行四边形面积为\(\|u\times v\|\)，平行六面体体积为\(|u\cdot(v\times w)|\)。

\begin{equation*}
 A=\begin{pmatrix}a&b\\c&d\end{pmatrix},\qquad \det A=ad-bc,\qquad
 A^{-1}=\frac1{ad-bc}\begin{pmatrix}d&-b\\-c&a\end{pmatrix}\quad(ad-bc\ne0).
\end{equation*}
平面仿射变换\(x=Au+t\)的面积倍率为\(|\det A|\)。

\section{空间几何体}

记\(B,P\)为底面积、底面周长，\(h\)为高；\(B_1,B_2\)与\(P_1,P_2\)分别对应两底面。\(l\)表示正棱锥、正棱台的斜高，或圆锥、圆台的母线长。表面积计入底面。
\begingroup
\renewcommand{\arraystretch}{1.15}
\begin{longtable}{@{}p{.22\linewidth}p{.31\linewidth}p{.41\linewidth}@{}}
\toprule
几何体 & 体积\(V\) & 表面积\(S\) \\
\midrule
\endhead
正方体（棱长\(a\)） & \(a^3\) & \(6a^2\) \\
长方体（\(a,b,c\)） & \(abc\) & \(2(ab+bc+ca)\) \\
正四面体（棱长\(a\)） & \(\sqrt2\,a^3/12\) & \(\sqrt3\,a^2\) \\
直棱柱 & \(Bh\) & \(2B+Ph\) \\
正棱锥 & \(Bh/3\) & \(B+Pl/2\) \\
正棱台 & \(\dfrac h3(B_1+B_2+\sqrt{B_1B_2})\) & \(B_1+B_2+\dfrac{P_1+P_2}{2}l\) \\
圆柱（半径\(r\)） & \(\pi r^2h\) & \(2\pi r(r+h)\) \\
圆锥（半径\(r\)） & \(\pi r^2h/3\) & \(\pi r(r+l)\) \\
圆台（半径\(R,r\)） & \(\dfrac{\pi h}{3}(R^2+Rr+r^2)\) & \(\pi(R^2+r^2)+\pi(R+r)l\) \\
球（半径\(r\)） & \(4\pi r^3/3\) & \(4\pi r^2\) \\
半球（半径\(r\)） & \(2\pi r^3/3\) & \(3\pi r^2\)（含圆形底面） \\
\bottomrule
\end{longtable}
\endgroup
圆锥的\(l=\sqrt{h^2+r^2}\)，圆台的\(l=\sqrt{h^2+(R-r)^2}\)。上述棱柱、棱锥、棱台的体积公式也适用于一般情形。

\section{复数}
\begin{equation*}
 \overline{a+bi}=a-bi,\qquad
 |z|^2=z\overline z,\qquad
 \frac1z=\frac{\overline z}{|z|^2}\quad(z\ne0).
\end{equation*}
对非零复数\(z=r(\cos\theta+i\sin\theta)=re^{i\theta}\)，
\begin{equation*}
 z^n=r^n(\cos n\theta+i\sin n\theta),\qquad n\in\mathbb Z.
\end{equation*}
若\(w=\rho e^{i\varphi}\ne0\)、\(n\in\mathbb N^*\)，方程\(z^n=w\)的全部根为
\begin{equation*}
 z_k=\rho^{1/n}e^{i(\varphi+2k\pi)/n},\qquad k=0,1,\ldots,n-1.
\end{equation*}
Euler公式为\(e^{i\theta}=\cos\theta+i\sin\theta\)。

\section{二次曲线}

以下\(a,b,p>0\)，切线公式中的\((x_0,y_0)\)在相应曲线上。
\begingroup
\renewcommand{\arraystretch}{1.3}
\begin{longtable}{@{}p{.41\linewidth}p{.53\linewidth}@{}}
\toprule
标准方程 & 切线方程 \\
\midrule
\endhead
\(\dfrac{x^2}{a^2}+\dfrac{y^2}{b^2}=1\) & \(\dfrac{xx_0}{a^2}+\dfrac{yy_0}{b^2}=1\) \\[6pt]
\(\dfrac{x^2}{a^2}-\dfrac{y^2}{b^2}=1\) & \(\dfrac{xx_0}{a^2}-\dfrac{yy_0}{b^2}=1\) \\[6pt]
\(y^2=2px\) & \(yy_0=p(x+x_0)\) \\
\bottomrule
\end{longtable}
\endgroup
椭圆的\(c^2=a^2-b^2\)（\(a\ge b\)），双曲线的\(c^2=a^2+b^2\)，两者离心率均为\(e=c/a\)。上述抛物线的焦点为\((p/2,0)\)，准线为\(x=-p/2\)。

椭圆可表示为\(x=a\cos t,y=b\sin t\)；双曲线可表示为\(x=\pm a\cosh t,y=b\sinh t\)；抛物线可表示为\(x=pt^2/2,y=pt\)。椭圆面积为\(\pi ab\)，上述双曲线的渐近线为\(y=\pm(b/a)x\)。

若非退化二次曲线在齐次坐标中写为\(X^TMX=0\)、\(M=M^T\)，点\(P\)的极线为\(P^TMX=0\)；切点\(P\)在曲线上时，这就是该点的切线。

\newpage
\section{求导公式}

下表中的参数按相应函数的定义取值，\(a>0,a\ne1\)，实数幂的基本定义域取\(x>0\)。
\begingroup
\renewcommand{\arraystretch}{1.1}
\begin{longtable}{@{}p{.44\linewidth}p{.48\linewidth}@{}}
\toprule
\(f(x)\) & \(f'(x)\) \\
\midrule
\endhead
\(x^\alpha\) & \(\alpha x^{\alpha-1}\) \\
\(e^x,\ a^x\) & \(e^x,\ a^x\ln a\) \\
\(\ln x,\ \log_a x\) & \(1/x,\ 1/(x\ln a)\) \\
\(\sin x,\ \cos x\) & \(\cos x,\ -\sin x\) \\
\(\tan x,\ \cot x\) & \(1/\cos^2x,\ -1/\sin^2x\) \\
\(\arcsin x,\ \arccos x\) & \(1/\sqrt{1-x^2},\ -1/\sqrt{1-x^2}\) \\
\(\arctan x\) & \(1/(1+x^2)\) \\
\(\sinh x,\ \cosh x\) & \(\cosh x,\ \sinh x\) \\
\(\tanh x,\ \coth x\) & \(1/\cosh^2x,\ -1/\sinh^2x\) \\
\bottomrule
\end{longtable}
\endgroup
\begin{equation*}
 \begin{aligned}
 (uv)'&=u'v+uv',&
 \left(\frac uv\right)'&=\frac{u'v-uv'}{v^2},\\
 (g\circ f)'(x)&=g'(f(x))f'(x),&
 (f^{-1})'(y)&=\frac1{f'(x)}\quad(y=f(x),\ f'(x)\ne0).
 \end{aligned}
\end{equation*}
对参数曲线\(x=x(t),y=y(t)\)，若\(x'(t)\ne0\)，则\(dy/dx=y'(t)/x'(t)\)。

\section{积分公式}

以下\(C\)为任意常数，\(\alpha\ne-1\)、\(a>0,a\ne1\)，各式在表达式有意义的区间上使用。
\begin{equation*}
 \begin{aligned}
 \int x^\alpha\,dx&=\frac{x^{\alpha+1}}{\alpha+1}+C,&
 \int\frac{dx}{x}&=\ln|x|+C,\\
 \int e^x\,dx&=e^x+C,&
 \int a^x\,dx&=\frac{a^x}{\ln a}+C,\\
 \int\sin x\,dx&=-\cos x+C,&
 \int\cos x\,dx&=\sin x+C,\\
 \int\frac{dx}{1+x^2}&=\arctan x+C,&
 \int\frac{dx}{\sqrt{1-x^2}}&=\arcsin x+C.
 \end{aligned}
\end{equation*}
在相应微积分定理的条件下，
\begin{equation*}
 \begin{gathered}
 \int_a^b f(x)\,dx=F(b)-F(a)\quad(F'=f),\qquad
 \int u\,dv=uv-\int v\,du,\\
 \int f(\varphi(t))\varphi'(t)\,dt=F(\varphi(t))+C,\\
 \frac d{dx}\int_{u(x)}^{v(x)}f(t)\,dt
 =f(v(x))v'(x)-f(u(x))u'(x).
 \end{gathered}
\end{equation*}

\section{Taylor公式与几何累积}

在相应Taylor定理的条件下，
\begin{equation*}
 f(x)=\sum_{k=0}^n\frac{f^{(k)}(a)}{k!}(x-a)^k+R_n(x),
\end{equation*}
其中常用的两种余项形式为
\begin{equation*}
 R_n(x)=\frac{f^{(n+1)}(\xi)}{(n+1)!}(x-a)^{n+1},\qquad
 R_n(x)=\frac1{n!}\int_a^x(x-t)^n f^{(n+1)}(t)\,dt.
\end{equation*}
第一式中的\(\xi\)位于\(a,x\)之间。二阶局部模型为\(f(a+h)=f(a)+f'(a)h+f''(a)h^2/2+o(h^2)\)。

对上下边界为\(f\ge g\)的区域、连续截面积为\(A(x)\)的几何体，以及光滑曲线，常用累积公式为
\begin{equation*}
 \begin{aligned}
 S&=\int_a^b(f-g)\,dx,&
 V&=\int_a^b A(x)\,dx,\\
 L&=\int_\alpha^\beta\sqrt{x'(t)^2+y'(t)^2}\,dt,&
 S_{\mathrm{rot}}&=2\pi\int_\alpha^\beta r(t)\sqrt{x'(t)^2+y'(t)^2}\,dt.
 \end{aligned}
\end{equation*}
\(r(t)\)为到旋转轴的距离；曲线或曲面按恰好覆盖一次取参数范围。曲率在速度非零处为
\begin{equation*}
 \kappa=\frac{|x'y''-y'x''|}{(x'^2+y'^2)^{3/2}}.
\end{equation*}

\section{微分方程}

一阶线性初值问题\(y'+py=q\)、\(y(t_0)=y_0\)中，记\(P(t)=\int_{t_0}^tp(s)\,ds\)，则
\begin{equation*}
 y(t)=e^{-P(t)}\left(y_0+\int_{t_0}^t e^{P(s)}q(s)\,ds\right).
\end{equation*}
二阶常系数齐次方程\(y''+ay'+by=0\)的特征方程为\(r^2+ar+b=0\)，通解按根型为
\begin{equation*}
 \begin{cases}
 C_1e^{r_1t}+C_2e^{r_2t},&r_1\ne r_2\text{为实根},\\
 (C_1+C_2t)e^{rt},&r\text{为重根},\\
 e^{\alpha t}(C_1\cos\beta t+C_2\sin\beta t),&\alpha\pm i\beta\ (\beta>0).
 \end{cases}
\end{equation*}
非齐次方程在已知一个特解\(y_p\)后，其通解为\(y_p+y_h\)，其中\(y_h\)是齐次通解。
\endgroup
